% Source: tex/chapters/ch06/04_ソフトウェアエンジニアリング運用の基礎.tex
\subsubsection{ソフトウェア導入}

ソフトウェア導入とは、ソフトウェアまたは更新版を利用者に使える状態にするため、対象環境に配置し、設定し、確認する作業である。導入では、旧版の削除、対象環境に合わせた設定、必要なディレクトリ、登録情報、環境変数の作成などが行われる。多くの場合、これらはスクリプトで実行される。

組込みシステムでは、電子的な配布だけでなく、物理媒体を使うこともある。いずれの場合でも、導入後には検証が必要である。導入が成功したか、必要なファイルが配置されたか、設定が正しいか、最低限の動作確認ができるかを確認する。
